\documentclass{article}
\usepackage{CJKutf8}
\usepackage{graphicx}
\usepackage{math_blog}
\usepackage{amsmath}
\usepackage{amssymb}
\usepackage{appendix}
\usepackage[normalem]{ulem}
\usepackage[thicklines]{cancel}
\usepackage{tikz}
\usetikzlibrary{positioning,arrows.meta}
\usepackage{pgfplots}

% Defensive: resolve possible \proof / theorem-env conflicts with math_blog
\makeatletter
\@ifundefined{proof}{}{\let\proof\@undefined \let\endproof\@undefined}
\@ifundefined{theorem}{\newtheorem{theorem}{Theorem}[section]}{}
\@ifundefined{lemma}{\newtheorem{lemma}[theorem]{Lemma}}{}
\@ifundefined{definition}{\newtheorem{definition}[theorem]{Definition}}{}
\@ifundefined{proposition}{\newtheorem{proposition}[theorem]{Proposition}}{}
\@ifundefined{corollary}{\newtheorem{corollary}[theorem]{Corollary}}{}
\@ifundefined{remark}{\newtheorem*{remark}{Remark}}{}
\@ifundefined{example}{\newtheorem*{example}{Example}}{}
\makeatother
\newenvironment{proof}[1][Proof]{%
  \par\smallskip\noindent\textit{\textbf{#1.}}\quad\ignorespaces
}{%
  \hfill$\square$\par\smallskip
}

\title{[SDE IV] Backward SDEs and the Nonlinear Feynman--Kac Formula}
\author{Anqiao Ouyang, Xuecheng Liu}
\date{Draft, \today}

\newcommand{\Title}{Backward SDEs and the Nonlinear Feynman--Kac Formula}
\newcommand{\Column}{Stochastic Differential Equations IV}
\newcommand{\Author}{Anqiao Ouyang, Xuecheng Liu}

% Math shortcuts
\newcommand{\E}{\mathbb{E}}
\newcommand{\Prob}{\mathbb{P}}
\newcommand{\R}{\mathbb{R}}
\newcommand{\F}{\mathcal{F}}
\newcommand{\Lc}{\mathcal{L}}
\newcommand{\Sc}{\mathcal{S}}
\newcommand{\Hc}{\mathcal{H}}

\begin{document}

\maketitle

\section{Introduction}

The previous three articles in this series built the forward picture of stochastic dynamics. In [SDE~I] we constructed Brownian motion and the surrounding measure-theoretic infrastructure. In [SDE~II] we developed the It\^o integral and the chain rule of stochastic calculus. In [SDE~III] we proved existence and uniqueness of solutions to forward SDEs and crossed over to PDEs via the Feynman--Kac formula: for the linear parabolic problem
\[
  \partial_t u + \Lc u = 0, \qquad u(T,x) = g(x),
\]
where $\Lc$ is the infinitesimal generator of $X$, the solution admits the probabilistic representation
\[
  u(t,x) = \E\big[g(X_T) \,\big|\, X_t = x\big].
\]

A natural question is what happens if we add a nonlinear source term to the PDE:
\[
  \partial_t u + \Lc u + f\!\big(t, x, u, \sigma^{\top}\nabla u\big) = 0, \qquad u(T, x) = g(x).
\]
The expectation in linear Feynman--Kac is taken at the fixed time $T$, and it does not see $u$ or $\nabla u$ inside the integrand. To bring those in, we need a representation that propagates information \emph{backward in time}, conditioning on the path of $X$ all along $[t, T]$. The object that does this is a \textbf{backward stochastic differential equation (BSDE)}.

A BSDE is, schematically, an SDE run from a terminal condition $Y_T = \xi$ back to time $0$:
\[
  -dY_t = f(t, Y_t, Z_t)\,dt - Z_t\, dB_t, \qquad Y_T = \xi.
\]
The unknown is the pair $(Y_t, Z_t)$, where $Y$ is a real-valued process and $Z$ is the integrand against Brownian motion. As we will see in \S2, $Z$ is forced upon us by adaptedness: it is not an extra data choice but a consequence of the martingale representation theorem. When $\xi = g(X_T)$ and $f$ depends on $(t, X_t, Y_t, Z_t)$ in a Markovian way, the BSDE solves precisely the semilinear PDE above, with the identification
\[
  Y_t = u(t, X_t), \qquad Z_t = \sigma^{\top}(t, X_t)\, \nabla u(t, X_t).
\]
This is the \textbf{nonlinear Feynman--Kac formula} of Pardoux and Peng.

The theory was initiated by Pardoux and Peng in 1990 and now sits at the heart of stochastic control, mathematical finance, and --- more recently --- deep-learning-based PDE solvers. In this article we will:

\begin{itemize}
  \item Define BSDEs precisely and explain why a pair $(Y, Z)$ is the right unknown (\S2).
  \item Prove the Pardoux--Peng existence and uniqueness theorem under Lipschitz conditions (\S3).
  \item Establish the comparison theorem and observe how it reflects the no-arbitrage structure of pricing (\S4).
  \item State and prove the nonlinear Feynman--Kac formula for Markovian BSDEs, side by side with the linear version from [SDE~III] (\S5).
  \item Apply the machinery to stochastic optimal control (HJB) and outline the Deep BSDE numerical method (\S6).
  \item Close with an outlook on forward--backward coupling, mean-field BSDEs, and rough paths (\S7).
\end{itemize}

\paragraph{Standing notation.} We work on a filtered probability space $(\Omega, \F, (\F_t)_{0 \le t \le T}, \Prob)$ carrying a $d$-dimensional Brownian motion $B = (B_t)_{0 \le t \le T}$, with $(\F_t)$ the augmented filtration of $B$ and a fixed horizon $T > 0$. Vectors are columns; $\sigma^{\top}$ is matrix transpose; $|\cdot|$ denotes Euclidean norm. ``Adapted'' means $(\F_t)$-adapted; ``predictable'' means $(\F_t)$-predictable. We will reuse $K$ as a generic constant whose value may change from line to line.

\section{Definition of BSDEs and the role of \texorpdfstring{$Z$}{Z}}

\subsection{Why a pair?}

Suppose we try to define a backward SDE in the naive way: pick a terminal condition $\xi \in L^2(\F_T)$ and a driver $f(t, y)$, and look for an $(\F_t)$-adapted process $Y$ satisfying
\[
  Y_T = \xi, \qquad -dY_t = f(t, Y_t)\,dt.
\]
This is just an ODE run backward, pathwise. The trouble is adaptedness: $Y_T = \xi$ is $\F_T$-measurable, so $Y_t$ defined by an ODE backward from $T$ inherits all the information in $\xi$, which violates $\F_t$-measurability at any earlier time $t < T$. In other words, knowing $Y_t$ would let us forecast $\xi$.

The remedy is to allow a martingale correction. By the martingale representation theorem (e.g.\ \cite{revuz-yor}, Ch.~V), every square-integrable Brownian martingale $M$ with $M_0 = \E[\xi]$ admits a unique predictable integrand $Z \in \Hc^2$ such that
\[
  M_t = M_0 + \int_0^t Z_s\, dB_s.
\]
Choosing $M_t := \E[\xi \mid \F_t]$ produces the $Z$ that ``carries'' the news about $\xi$ revealed up to time $t$. The BSDE is the integral identity that combines the deterministic driver $f$ with this Brownian correction $Z\, dB$.

\begin{definition}[BSDE solution]\label{def:bsde}
Let $\xi \in L^2(\F_T; \R)$ and $f : [0,T] \times \Omega \times \R \times \R^d \to \R$ be progressively measurable. A pair of processes $(Y, Z)$ --- with $Y$ continuous, real-valued, $(\F_t)$-adapted, and $Z$ $(\F_t)$-predictable, $\R^d$-valued --- is called a \emph{solution} to the BSDE $(\xi, f)$ if
\begin{equation}\label{eq:bsde}
  Y_t = \xi + \int_t^T f(s, Y_s, Z_s)\,ds - \int_t^T Z_s\, dB_s, \qquad t \in [0, T],
\end{equation}
$\Prob$-a.s., and additionally
\[
  \E\!\Big[\sup_{0 \le t \le T}|Y_t|^2\Big] < \infty, \qquad \E\!\int_0^T |Z_s|^2\,ds < \infty.
\]
\end{definition}

We give the solution spaces names:
\[
  \Sc^2 := \Big\{Y \text{ adapted, continuous} : \|Y\|_{\Sc^2}^2 := \E\!\big[\sup_t |Y_t|^2\big] < \infty\Big\},
\]
\[
  \Hc^2 := \Big\{Z \text{ predictable} : \|Z\|_{\Hc^2}^2 := \E\!\int_0^T |Z_s|^2\,ds < \infty\Big\}.
\]
Both are Banach (in fact Hilbert), and the equation \eqref{eq:bsde} is to be read in $\Sc^2 \times \Hc^2$.

\begin{remark}
Differentiating \eqref{eq:bsde} gives the differential form
\[
  -dY_t = f(t, Y_t, Z_t)\,dt - Z_t\, dB_t, \qquad Y_T = \xi.
\]
The minus sign in front of $dY$ is bookkeeping: it records that time is being run backward from $T$. Some authors absorb it into the driver by writing $\widetilde{f} = -f$.
\end{remark}

\begin{figure}[h]
\centering
\begin{tikzpicture}[scale=1.0]
\draw[->] (-0.4,0) -- (10.7,0) node[right] {\small time};
\node[below] at (0,0) {\small$0$};
\node[below] at (3.0,-0.05) {\small$t$};
\node[below] at (9.5,-0.05) {\small$T$};
\foreach \x in {0,3.0,9.5} \draw (\x,0.08) -- (\x,-0.08);

% Forward X path
\draw[blue, thick] plot[smooth] coordinates {
  (3.0,1.6) (3.4,1.5) (3.8,1.85) (4.2,1.7) (4.6,2.1) (5.0,1.9)
  (5.4,2.3) (5.8,2.6) (6.2,2.4) (6.6,2.85) (7.0,3.05) (7.4,2.9)
  (7.8,3.3) (8.2,3.45) (8.6,3.25) (9.0,3.55) (9.5,3.4)};
\fill[blue] (3.0,1.6) circle (2pt) node[left,blue] {\small$X^{t,x}_t=x$};
\fill[blue] (9.5,3.4) circle (2pt) node[above right,blue] {\small$X^{t,x}_T$};
\node[blue] at (5.4,3.55) {\small forward SDE: $X^{t,x}$};
\draw[->, blue!70, thick] (3.4,1.05) -- (5.6,1.05);
\node[blue!70, below, font=\scriptsize] at (4.5,1.05) {time $\longrightarrow$};

% Backward Y path
\draw[red, thick] plot[smooth] coordinates {
  (3.0,-1.4) (3.4,-1.42) (3.8,-1.32) (4.2,-1.5) (4.6,-1.4) (5.0,-1.55)
  (5.4,-1.45) (5.8,-1.6) (6.2,-1.55) (6.6,-1.7) (7.0,-1.62) (7.4,-1.78)
  (7.8,-1.7) (8.2,-1.85) (8.6,-1.78) (9.0,-1.92) (9.5,-2.0)};
\fill[red] (3.0,-1.4) circle (2pt) node[left,red] {\small$Y_t=u(t,x)$};
\fill[red] (9.5,-2.0) circle (2pt) node[right,red] {\small$\xi=g(X_T)$};
\node[red] at (5.7,-2.55) {\small backward SDE: $(Y,Z)$};
\draw[<-, red!70, thick] (5.7,-0.55) -- (7.9,-0.55);
\node[red!70, above, font=\scriptsize] at (6.8,-0.55) {$\longleftarrow$ time};

\draw[gray!50, dashed] (3.0,1.6) -- (3.0,-1.4);
\draw[gray!50, dashed] (9.5,3.4) -- (9.5,-2.0);

\node[align=center,font=\scriptsize, gray!60!black] at (10.6,0.85) {same $dB_s$};
\draw[->, gray!60!black] (10.3,0.65) -- (9.7,0.3);
\draw[->, gray!60!black] (10.3,0.65) -- (9.7,-1.0);
\end{tikzpicture}
\caption{The forward--backward picture underlying the nonlinear Feynman--Kac formula. The diffusion $X^{t,x}$ is launched from $x$ at time $t$ and runs \emph{forward} to a random terminal value $X_T$; the terminal payoff $\xi = g(X_T)$ is then transported \emph{backward} by the BSDE to produce $Y_t = u(t,x)$. Both processes are driven by the same Brownian increments $dB_s$ on $[t,T]$: the forward equation uses them through the diffusion coefficient $\sigma$, the backward equation through the integrand $Z$. Adaptedness forces $Y_t$ to be a function of the history up to $t$ only; this is what makes $Z$ unavoidable.}
\label{fig:fwd-bwd}
\end{figure}

\subsection{Two warm-up examples}

\begin{example}[$f \equiv 0$: linear Feynman--Kac in disguise]
The BSDE $(\xi, 0)$ reduces to $Y_t = \xi - \int_t^T Z_s\,dB_s$. Taking conditional expectations,
\[
  Y_t = \E[\xi \mid \F_t],
\]
and $Z$ is the integrand from the martingale representation of $M_t = \E[\xi \mid \F_t]$. The BSDE with zero driver is exactly the linear Feynman--Kac formula. Specializing to $\xi = g(X_T)$ recovers [SDE~III, Thm.~4.1].
\end{example}

\begin{example}[Linear driver and explicit solution]
For $f(t, y, z) = \alpha_t\, y + \beta_t \cdot z + \gamma_t$ with bounded adapted coefficients, define the discount and change-of-measure factors
\[
  \Gamma_t := \exp\!\Big(\int_0^t \alpha_s\, ds\Big), \qquad \frac{d\Prob^\beta}{d\Prob}\bigg|_{\F_T} := \exp\!\Big(\int_0^T \beta_s\, dB_s - \tfrac12\!\int_0^T |\beta_s|^2 ds\Big).
\]
Then
\[
  \Gamma_t\, Y_t = \E^{\Prob^\beta}\!\Big[ \Gamma_T\, \xi + \int_t^T \Gamma_s\, \gamma_s\, ds \,\Big|\, \F_t\Big].
\]
This formula is the BSDE analogue of variation of parameters; we will use it inside the proof of the comparison theorem.
\end{example}

\section{Existence and Uniqueness}

The cornerstone result of the linear theory is the following.

\begin{theorem}[Pardoux--Peng, 1990]\label{thm:pp}
Assume
\begin{enumerate}
  \item[(i)] $\xi \in L^2(\F_T)$;
  \item[(ii)] $(t, \omega) \mapsto f(t, \omega, 0, 0)$ belongs to $\Hc^2$;
  \item[(iii)] $f$ is uniformly Lipschitz in $(y, z)$: there exists $K \ge 0$ with
    \[
      |f(t, \omega, y, z) - f(t, \omega, y', z')| \le K\big(|y - y'| + |z - z'|\big)
    \]
    $dt \otimes d\Prob$-a.e.
\end{enumerate}
Then the BSDE \eqref{eq:bsde} admits a unique solution $(Y, Z) \in \Sc^2 \times \Hc^2$.
\end{theorem}

The proof rests on a contraction argument in the spirit of Picard iteration for forward SDEs ([SDE~III, \S2]), with one new ingredient: the contraction happens in a \emph{weighted} norm whose exponential weight absorbs the Lipschitz constant.

\subsection{An a priori estimate}

\begin{lemma}[A priori bound]\label{lem:apriori}
If $(Y, Z) \in \Sc^2 \times \Hc^2$ solves \eqref{eq:bsde}, then for every $\beta > 2K + 2K^2$,
\[
  \E\!\int_0^T e^{\beta s}\big(|Y_s|^2 + |Z_s|^2\big)\,ds \le C \!\left( \E\big[ e^{\beta T}|\xi|^2 \big] + \E\!\int_0^T e^{\beta s}|f(s, 0, 0)|^2\, ds \right),
\]
for a constant $C$ depending only on $K$ and $\beta$.
\end{lemma}

\begin{proof}[Sketch]
Apply It\^o's formula to $\varphi_t := e^{\beta t} |Y_t|^2$:
\[
  d\varphi_t = \beta e^{\beta t}|Y_t|^2\,dt + 2 e^{\beta t} Y_t\,dY_t + e^{\beta t} |Z_t|^2\,dt.
\]
Substituting $dY_t = -f(t, Y_t, Z_t)\,dt + Z_t\,dB_t$ and integrating from $t$ to $T$,
\begin{align*}
  e^{\beta t}|Y_t|^2 + \int_t^T e^{\beta s}\big(\beta |Y_s|^2 + |Z_s|^2\big)\,ds
  &= e^{\beta T}|\xi|^2 + 2\int_t^T e^{\beta s}\, Y_s\, f(s, Y_s, Z_s)\,ds \\
  &\quad - 2\int_t^T e^{\beta s}\, Y_s\, Z_s\, dB_s.
\end{align*}
Take expectations (the It\^o integral has zero mean given local boundedness, which one bootstraps from $(Y, Z) \in \Sc^2 \times \Hc^2$). Decompose
\[
  Y_s\, f(s, Y_s, Z_s) = Y_s\, f(s, 0, 0) + Y_s\big(f(s, Y_s, Z_s) - f(s, 0, 0)\big),
\]
use Lipschitz on the second piece and the elementary $2|y||z| \le \varepsilon|y|^2 + \varepsilon^{-1}|z|^2$, then choose $\varepsilon$ small and $\beta$ large enough to absorb $\beta|Y|^2$ and $|Z|^2$ on the left. The constant $\beta > 2K + 2K^2$ does the job.
\end{proof}

A consequence (apply the lemma with $\xi = 0$ and $f \equiv 0$, to two solutions): \emph{uniqueness}. Two solutions $(Y, Z)$, $(Y', Z')$ to the same BSDE have difference $(Y - Y', Z - Z')$ solving a BSDE with zero data; the estimate forces both to vanish.

\subsection{Existence by contraction}

Define a map $\Phi : \Sc^2 \times \Hc^2 \to \Sc^2 \times \Hc^2$ as follows. Given $(y, z)$, set $(Y, Z) := \Phi(y, z)$ to be the solution of the BSDE with \emph{frozen} driver $f^{y,z}(t, \omega) := f(t, \omega, y_t, z_t)$ and terminal condition $\xi$:
\[
  Y_t = \xi + \int_t^T f(s, y_s, z_s)\,ds - \int_t^T Z_s\, dB_s.
\]
For this linear sub-problem (no $Y$- or $Z$-dependence on the right) the solution can be written explicitly: set
\[
  N_t := \E\!\left[\xi + \int_0^T f(s, y_s, z_s)\,ds \,\bigg|\, \F_t\right].
\]
Then $N$ is a square-integrable martingale; by Brownian representation, $N_t = N_0 + \int_0^t Z_s\, dB_s$ for a unique predictable $Z \in \Hc^2$. Put
\[
  Y_t := N_t - \int_0^t f(s, y_s, z_s)\,ds,
\]
and verify directly that $(Y, Z)$ satisfies the sub-problem and lies in $\Sc^2 \times \Hc^2$.

The next lemma turns the a priori bound into a quantitative Lipschitz estimate on $\Phi$.

\begin{lemma}[Difference estimate]\label{lem:diff}
For $(y, z), (y', z') \in \Sc^2 \times \Hc^2$, let $(Y, Z) := \Phi(y, z)$, $(Y', Z') := \Phi(y', z')$. For every $\beta \ge 2$,
\[
  \|(Y - Y', Z - Z')\|_\beta^2 \;\le\; \frac{4K^2}{\beta}\, \|(y - y', z - z')\|_\beta^2,
\]
where $\|(U, V)\|_\beta^2 := \E\!\int_0^T e^{\beta s}\big(|U_s|^2 + |V_s|^2\big)\,ds$ and $K$ is the Lipschitz constant of $f$ in $(y, z)$.
\end{lemma}

\begin{proof}
Write $\Delta y := y - y'$, $\Delta z := z - z'$, $\Delta Y := Y - Y'$, $\Delta Z := Z - Z'$, and $\Delta f_s := f(s, y_s, z_s) - f(s, y'_s, z'_s)$. The two frozen BSDEs are identical at time $T$ ($\Delta Y_T = 0$) and yield
\[
  d\Delta Y_t = -\Delta f_t\, dt + \Delta Z_t\, dB_t.
\]
Apply It\^o's formula to $e^{\beta s} |\Delta Y_s|^2$ on $[0, T]$, take expectations, and use $\Delta Y_T = 0$:
\[
  \E\!\int_0^T e^{\beta s}\!\big(\beta |\Delta Y_s|^2 + |\Delta Z_s|^2\big)\, ds \;=\; 2\,\E\!\int_0^T e^{\beta s}\, \Delta Y_s\, \Delta f_s\, ds.
\]
By Lipschitz, $|\Delta f_s| \le K(|\Delta y_s| + |\Delta z_s|)$. Applying $2ab \le \tfrac{\beta}{4}a^2 + \tfrac{4}{\beta} b^2$ to each cross term,
\[
  2\, \Delta Y_s\, \Delta f_s \;\le\; \tfrac{\beta}{2}|\Delta Y_s|^2 + \tfrac{4 K^2}{\beta}\!\big(|\Delta y_s|^2 + |\Delta z_s|^2\big).
\]
Absorb the $\tfrac{\beta}{2}|\Delta Y|^2$ on the left and use $\beta \ge 2$ to bound the $|\Delta Z|^2$ coefficient below by $1$, giving
\[
  \E\!\int_0^T e^{\beta s}\big(|\Delta Y_s|^2 + |\Delta Z_s|^2\big)\, ds \;\le\; \frac{4K^2}{\beta}\, \E\!\int_0^T e^{\beta s}\big(|\Delta y_s|^2 + |\Delta z_s|^2\big)\, ds.
\]
\end{proof}

\noindent Choosing any $\beta > 4K^2$ (e.g.\ $\beta = 8 K^2 + 2$) makes $\Phi$ a strict contraction in $\|\cdot\|_\beta$ with ratio $\sqrt{4K^2/\beta} < 1$. Banach's fixed-point theorem on the complete metric space $(\Sc^2 \times \Hc^2, \|\cdot\|_\beta)$ produces a unique fixed point, which by construction solves the BSDE \eqref{eq:bsde}. Combined with the uniqueness corollary, this completes the proof of Theorem~\ref{thm:pp}. $\square$

\begin{remark}[Why a weighted norm?]
The unweighted $L^2$-norm fails because the contraction constant from the same calculation reads $4 K^2 T$, which is $\ge 1$ once $T$ is moderately large. The exponential weight $e^{\beta s}$ rescales time so that one can spend the full Lipschitz budget locally without the global horizon $T$ entering. The same trick (under the name ``Bielecki norm'') is the standard cure for Picard-type contractions on long intervals --- it appeared in [SDE~III, \S2] for forward SDEs and is even more crucial here, because the BSDE cannot be solved by short-step concatenation: shrinking $T$ does not help, since the terminal datum lives at $T$.
\end{remark}

\begin{remark}[What if $f$ is only monotone in $y$?]
Lipschitz can be replaced by ``monotonicity in $y$ plus polynomial growth in $y$, Lipschitz in $z$'' (Pardoux 1999). The proof becomes more delicate but follows the same outline. Quadratic growth in $z$ (Kobylanski 2000) requires bounded $\xi$ and entirely different techniques (exponential transforms).
\end{remark}

\section{The Comparison Theorem}

A defining feature of BSDEs --- and one with no counterpart in the forward theory --- is the following monotonicity in the data.

\begin{theorem}[Comparison]\label{thm:comparison}
Let $(Y^1, Z^1)$ and $(Y^2, Z^2)$ solve the BSDEs $(\xi^1, f^1)$ and $(\xi^2, f^2)$ respectively, with both drivers Lipschitz. Assume
\[
  \xi^1 \le \xi^2 \text{ a.s.}, \qquad f^1(t, Y^1_t, Z^1_t) \le f^2(t, Y^1_t, Z^1_t) \quad dt \otimes d\Prob\text{-a.e.}
\]
Then $Y^1_t \le Y^2_t$ a.s.\ for every $t \in [0, T]$.
\end{theorem}

\begin{proof}[Sketch]
Set $\Delta Y := Y^2 - Y^1$ and $\Delta Z := Z^2 - Z^1$. Subtracting the two BSDEs,
\[
  -d\Delta Y_t = \big[f^2(t, Y^2, Z^2) - f^1(t, Y^1, Z^1)\big]\,dt - \Delta Z_t\, dB_t.
\]
Linearize the driver difference:
\[
  f^2(t, Y^2_t, Z^2_t) - f^1(t, Y^1_t, Z^1_t) = \alpha_t \Delta Y_t + \beta_t \cdot \Delta Z_t + \gamma_t,
\]
where $\alpha, \beta$ are bounded (by Lipschitz) and $\gamma_t := f^2(t, Y^1_t, Z^1_t) - f^1(t, Y^1_t, Z^1_t) \ge 0$ by hypothesis. So $\Delta Y$ solves a \emph{linear} BSDE with positive source. By a change of measure (Girsanov, [SDE~III, \S5]) under
\[
  \frac{d\Prob^\beta}{d\Prob}\bigg|_{\F_T} = \exp\!\Big(\int_0^T \beta_s\, dB_s - \tfrac12\!\int_0^T |\beta_s|^2\, ds\Big),
\]
the term $\beta \cdot \Delta Z$ is absorbed into the Brownian motion. Applying the linear-driver formula from \S2.2,
\[
  \Delta Y_t = \E^{\Prob^\beta}\!\left[ e^{\int_t^T \alpha_s\, ds}\, \Delta \xi + \int_t^T e^{\int_t^s \alpha_r\, dr}\, \gamma_s\, ds \,\bigg|\, \F_t \right].
\]
The integrand is a sum of two nonnegative quantities ($\Delta \xi \ge 0$ and $\gamma \ge 0$), so $\Delta Y_t \ge 0$.
\end{proof}

\begin{remark}[Finance interpretation]
In an option-pricing model, $Y_t$ represents the price at time $t$ of a contract paying $\xi$ at maturity $T$, and $f$ encodes the cost of hedging (interest rates, dividends, portfolio constraints). The comparison theorem says: \emph{a more valuable payoff or a cheaper hedge yields a higher price.} This is the dynamic-programming form of no-arbitrage; it is what powers, for instance, the upper-bound pricing under model uncertainty.
\end{remark}

\begin{corollary}[Strict comparison]\label{cor:strict-comp}
Under the hypotheses of Theorem~\ref{thm:comparison}, if additionally
\[
  \Prob(\xi^1 < \xi^2) > 0 \qquad\text{or}\qquad f^1 < f^2 \text{ at } (Y^1, Z^1)\text{ on a set of positive } dt\otimes d\Prob\text{-measure,}
\]
then $Y^1_0 < Y^2_0$.
\end{corollary}

\begin{proof}
In the representation derived in the proof of Theorem~\ref{thm:comparison},
\[
  \Delta Y_0 = \E^{\Prob^\beta}\!\left[ e^{\int_0^T \alpha_s\, ds}\, \Delta \xi + \int_0^T e^{\int_0^s \alpha_r\, dr}\, \gamma_s\, ds \right],
\]
the integrand is nonnegative and, under one of the two hypotheses, strictly positive on a set of positive $\Prob^\beta$-measure (which coincides with positive $\Prob$-measure by absolute continuity of $\Prob^\beta$). Hence $\Delta Y_0 > 0$.
\end{proof}

\subsection{$g$-expectations and dynamic risk measures}

The comparison theorem is what makes the following construction well-defined. Fix a Lipschitz driver $g : [0,T] \times \R \times \R^d \to \R$ satisfying $g(t, y, 0) = 0$, and for $\xi \in L^2(\F_T)$ let $(Y^\xi, Z^\xi)$ be the unique solution of the BSDE $(\xi, g)$. Define the \emph{$g$-expectation} and its conditional version by
\[
  \mathcal{E}_g[\xi] := Y^\xi_0, \qquad \mathcal{E}_g[\xi \mid \F_t] := Y^\xi_t.
\]
The condition $g(\cdot, \cdot, 0) = 0$ ensures that $\mathcal{E}_g[c] = c$ for constants. The operator $\mathcal{E}_g$ is in general nonlinear in $\xi$ (linearity holds iff $g$ is linear in $(y, z)$), but it retains the structural features of an expectation:
\begin{enumerate}
  \item[(M)] \emph{Monotonicity:} $\xi^1 \le \xi^2 \Rightarrow \mathcal{E}_g[\xi^1 \mid \F_t] \le \mathcal{E}_g[\xi^2 \mid \F_t]$ (Theorem~\ref{thm:comparison}).
  \item[(T)] \emph{Time-consistency:} $\mathcal{E}_g[\mathcal{E}_g[\xi \mid \F_s] \mid \F_t] = \mathcal{E}_g[\xi \mid \F_t]$ for $t \le s$ (flow property of the BSDE).
  \item[(L)] \emph{Locality / zero-one law:} $\mathbf{1}_A \mathcal{E}_g[\xi \mid \F_t] = \mathbf{1}_A \mathcal{E}_g[\mathbf{1}_A \xi \mid \F_t]$ for $A \in \F_t$.
\end{enumerate}
When $g$ is positively homogeneous and subadditive in $(y, z)$, $\rho_t(\xi) := \mathcal{E}_g[-\xi \mid \F_t]$ is a \emph{dynamic coherent risk measure} in the sense of Artzner--Delbaen--Eber--Heath; convex $g$ gives a dynamic convex risk measure. Conversely, a theorem of Coquet--Hu--M\'emin--Peng characterizes which dynamically consistent nonlinear expectations arise this way: precisely those dominated by a $g_\mu$-expectation for the linear driver $g_\mu(z) = \mu|z|$. Nonlinearity in $g$ encodes hedging frictions (bid--ask spreads, collateral, capital charges) and Knightian uncertainty about the law of $X$.

\section{The Nonlinear Feynman--Kac Formula}

We now specialize to the Markovian setting that makes contact with PDEs.

\subsection{Setting}

Fix $(t, x) \in [0, T) \times \R^n$. Consider the forward SDE
\begin{equation}\label{eq:fwd}
  dX^{t,x}_s = b(s, X^{t,x}_s)\,ds + \sigma(s, X^{t,x}_s)\, dB_s, \qquad X^{t,x}_t = x,\ s \in [t, T],
\end{equation}
with $b : [0, T] \times \R^n \to \R^n$ and $\sigma : [0, T] \times \R^n \to \R^{n \times d}$ satisfying the standing Lipschitz and linear-growth hypotheses from [SDE~III, \S2]. Coupled to it, consider the BSDE
\begin{equation}\label{eq:fkbsde}
  Y^{t,x}_s = g(X^{t,x}_T) + \int_s^T f\!\big(r, X^{t,x}_r, Y^{t,x}_r, Z^{t,x}_r\big)\,dr - \int_s^T Z^{t,x}_r\, dB_r,
\end{equation}
with $g : \R^n \to \R$ continuous of polynomial growth and $f : [0, T] \times \R^n \times \R \times \R^d \to \R$ continuous and uniformly Lipschitz in $(y, z)$. By Theorem~\ref{thm:pp}, the BSDE \eqref{eq:fkbsde} has a unique solution.

Define the candidate PDE solution
\[
  u(t, x) := Y^{t,x}_t.
\]
Note that $Y^{t,x}_t$ is $\F_t$-measurable, and because $X^{t,x}_t = x$ deterministically, the value $Y^{t,x}_t$ is actually deterministic in $(t, x)$ --- so $u(t, x)$ is a number, not a random variable.

Recall the generator of $X$:
\[
  \Lc \varphi(t, x) := b(t, x) \cdot \nabla \varphi(t, x) + \tfrac12\, \mathrm{tr}\!\big(\sigma\sigma^{\top}(t, x)\, \nabla^2 \varphi(t, x)\big).
\]

\begin{theorem}[Nonlinear Feynman--Kac]\label{thm:nlfk}
\hfill
\begin{enumerate}
  \item[(a)] (Verification) If a function $v \in C^{1,2}([0, T] \times \R^n)$ of polynomial growth solves the semilinear parabolic problem
  \begin{equation}\label{eq:semilinpde}
    \partial_t v(t, x) + \Lc v(t, x) + f\!\big(t, x, v(t, x), \sigma^{\top}(t, x)\, \nabla v(t, x)\big) = 0,
  \end{equation}
  with terminal condition $v(T, x) = g(x)$, then for every $(t, x)$,
  \[
    Y^{t,x}_s = v(s, X^{t,x}_s), \qquad Z^{t,x}_s = \sigma^{\top}(s, X^{t,x}_s)\, \nabla v(s, X^{t,x}_s),
  \]
  is the unique solution to the BSDE \eqref{eq:fkbsde}. In particular $u(t, x) = v(t, x)$.

  \item[(b)] (Representation) Conversely, if the candidate $u(t, x) := Y^{t,x}_t$ defined from the BSDE is in $C^{1,2}$, then $u$ solves \eqref{eq:semilinpde} with terminal condition $g$.
\end{enumerate}
\end{theorem}

\begin{proof}[Verification half]
Suppose $v \in C^{1,2}$ solves \eqref{eq:semilinpde}. Apply It\^o's formula to $v(s, X^{t,x}_s)$:
\[
  dv(s, X_s) = \big(\partial_s v + \Lc v\big)(s, X_s)\,ds + \nabla v(s, X_s)^{\top} \sigma(s, X_s)\, dB_s.
\]
By the PDE \eqref{eq:semilinpde}, the drift becomes $-f(s, X_s, v, \sigma^{\top} \nabla v)$, so
\[
  -dv(s, X_s) = f\!\big(s, X_s, v(s, X_s), \sigma^{\top}(s, X_s) \nabla v(s, X_s)\big)\,ds - \sigma^{\top}(s, X_s) \nabla v(s, X_s)\, dB_s,
\]
with terminal value $v(T, X_T) = g(X_T)$. This is exactly the BSDE \eqref{eq:fkbsde} with $Y_s := v(s, X_s)$ and $Z_s := \sigma^{\top}(s, X_s) \nabla v(s, X_s)$. By the uniqueness half of Theorem~\ref{thm:pp}, this is the only solution.
\end{proof}

The converse (representation) half is more subtle: showing $u \in C^{1,2}$ requires either Malliavin smoothness of the flow $x \mapsto X^{t,x}_T$, or a regularization-by-uniformly-elliptic-noise argument. We refer to \cite{el-karoui-quenez} and \cite{pardoux-rascanu} for two complementary approaches.

\begin{remark}[Why $\sigma^{\top} \nabla u$ and not $\nabla u$?]
The forward SDE drives the noise through $\sigma$. It\^o's formula on $v(s, X_s)$ produces $(\nabla v)^{\top} \sigma\, dB$, so the integrand against $dB$ --- which is precisely what $Z$ is --- equals $\sigma^{\top} \nabla v$. When $\sigma = I_d$ (driftless Brownian noise on $\R^d$), this collapses to the familiar $Z = \nabla u$.
\end{remark}

\begin{remark}[The Malliavin viewpoint]\label{rem:malliavin}
If $\xi$ and the data are Malliavin-differentiable in the sense of Nualart, then so is the solution $(Y, Z)$, and one has the identification
\begin{equation}\label{eq:Z-DY}
  Z_t \;=\; D_t Y_t \qquad dt \otimes d\Prob\text{-a.e.,}
\end{equation}
where $D_t$ is the Malliavin derivative (Pardoux--Peng 1992; El Karoui--Peng--Quenez 1997). The pair $(D_s Y_t, D_s Z_t)_{s \le t}$ itself solves a linear BSDE obtained by differentiating \eqref{eq:bsde} formally, which feeds back the regularity of the data into regularity of the solution. In the Markovian setting the chain rule and the tangent-flow formula $D_s X_t = \mathbf{1}_{s \le t}\,\nabla_x X^{s,X_s}_t\, \sigma(s, X_s)$ collapse \eqref{eq:Z-DY} to
\[
  Z_t = \sigma^{\top}(t, X_t)\, \nabla u(t, X_t),
\]
matching part (a) of Theorem~\ref{thm:nlfk} \emph{without} assuming $u \in C^{1,2}$ a priori. The Malliavin route is in fact how the converse half is proved: one establishes Sobolev (rather than classical) regularity of $u$ on the support of the law of $X$, then upgrades to $C^{1,2}$ via parabolic Schauder estimates when $\sigma\sigma^{\top}$ is uniformly elliptic.
\end{remark}

\subsection{Linear vs.\ nonlinear: a side-by-side comparison}

\begin{center}
\begin{tabular}{p{0.45\textwidth}|p{0.45\textwidth}}
\textbf{Linear Feynman--Kac (\,[SDE~III])} & \textbf{Nonlinear Feynman--Kac (this article)} \\
\hline
PDE: $\partial_t u + \Lc u = 0$, $u(T, x) = g(x)$ &
PDE: $\partial_t u + \Lc u + f(t, x, u, \sigma^{\top} \nabla u) = 0$, $u(T, x) = g(x)$ \\
\hline
Probabilistic side: forward SDE only ($X^{t,x}$) &
Probabilistic side: forward SDE \emph{plus} BSDE for $(Y, Z)$ \\
\hline
Representation: $u(t, x) = \E[g(X^{t,x}_T)]$ &
Representation: $u(t, x) = Y^{t,x}_t$, the initial value of the BSDE \\
\hline
Gradient $\nabla u$: not represented; need Bismut--Elworthy or Malliavin &
Gradient: $Z^{t,x}_t = \sigma^{\top}(t, x)\, \nabla u(t, x)$ --- read off directly from the BSDE \\
\hline
PDE drift has no $u$-dependence &
$f$ may depend on $u$ and $\nabla u$ in any (Lipschitz) nonlinear way \\
\end{tabular}
\end{center}

The right column reduces to the left when $f \equiv 0$, which is the warm-up example from \S2.2.

\subsection{A worked example: a Fisher--KPP-type equation}

Consider on $\R$ the semilinear equation
\[
  \partial_t u + \tfrac12 \partial_{xx} u + u(1 - u) = 0, \qquad u(T, x) = \mathbf{1}_{x > 0},
\]
with $b \equiv 0$, $\sigma \equiv 1$. The associated forward process is Brownian motion $X_s = x + B_s - B_t$, and the BSDE is
\[
  Y_s = \mathbf{1}_{x + B_T - B_t > 0} + \int_s^T Y_r(1 - Y_r)\, dr - \int_s^T Z_r\, dB_r.
\]
Theorem~\ref{thm:nlfk} identifies $u(t, x) = Y^{t,x}_t$. Note that $u \in [0, 1]$ pathwise --- one can check that the driver $y \mapsto y(1 - y)$ preserves the interval $[0, 1]$ via the comparison theorem applied to the constant solutions $y \equiv 0$ and $y \equiv 1$. This is exactly the maximum principle for the Fisher--KPP equation, recovered probabilistically.

The BSDE route also gives a Monte Carlo algorithm for $u$ that never touches a grid: simulate $X$ forward, then march $(Y, Z)$ backward by a discretization. This is the entry point to the Deep BSDE method of \S6.2.

\begin{figure}[h]
\centering
\begin{tikzpicture}
\begin{axis}[
    width=0.85\textwidth, height=6cm,
    xlabel={$s$}, ylabel={$Y_s$},
    xmin=0, xmax=1.05,
    ymin=-0.08, ymax=1.12,
    no markers,
    samples=100, smooth,
]
\addplot[gray, dashed, very thick, domain=0:1] {1};
\addplot[gray, dashed, very thick, domain=0:1] {0};
\node[gray, anchor=west, font=\scriptsize] at (axis cs: 1.005, 1) {$y\equiv1$};
\node[gray, anchor=west, font=\scriptsize] at (axis cs: 1.005, 0) {$y\equiv0$};

\addplot[blue, thick] coordinates {
(0,0.66)(0.1,0.69)(0.2,0.72)(0.3,0.75)(0.4,0.79)(0.5,0.82)(0.6,0.85)(0.7,0.89)(0.8,0.92)(0.9,0.96)(1.0,1.0)
};
\addplot[red, thick] coordinates {
(0,0.32)(0.1,0.27)(0.2,0.23)(0.3,0.21)(0.4,0.17)(0.5,0.13)(0.6,0.09)(0.7,0.07)(0.8,0.04)(0.9,0.02)(1.0,0)
};
\addplot[green!55!black, thick] coordinates {
(0,0.49)(0.1,0.53)(0.2,0.50)(0.3,0.58)(0.4,0.62)(0.5,0.66)(0.6,0.70)(0.7,0.78)(0.8,0.86)(0.9,0.94)(1.0,1.0)
};
\addplot[orange, thick] coordinates {
(0,0.51)(0.1,0.45)(0.2,0.38)(0.3,0.33)(0.4,0.28)(0.5,0.22)(0.6,0.16)(0.7,0.11)(0.8,0.06)(0.9,0.02)(1.0,0)
};
\end{axis}
\end{tikzpicture}
\caption{Four simulated $Y$-trajectories for the Fisher--KPP BSDE on $[0, T] = [0, 1]$ with $x = 0$, generated by Euler--Maruyama discretization of the forward Brownian path together with an explicit backward sweep for $(Y, Z)$. All trajectories remain in $[0, 1]$: the constants $y \equiv 0$ and $y \equiv 1$ are themselves BSDE solutions for the boundary terminal data $\xi = 0$ and $\xi = 1$, and the comparison theorem (Theorem~\ref{thm:comparison}) sandwiches every solution between them. The terminal values realize the indicator $\mathbf{1}_{B_T > 0}$ along each Brownian path; the initial values $Y_0$ are Monte Carlo estimates of $u(0, 0) \approx \tfrac12$.}
\label{fig:fisher-kpp}
\end{figure}

\section{Applications}

\subsection{Stochastic optimal control and HJB}

Consider a controlled diffusion
\[
  dX^\alpha_s = b(s, X^\alpha_s, \alpha_s)\, ds + \sigma(s, X^\alpha_s)\, dB_s,
\]
where the control $\alpha$ takes values in a compact action set $A \subset \R^k$ and is required to be adapted. With running cost $\ell$ and terminal cost $g$, define the value function
\[
  V(t, x) := \inf_{\alpha}\, \E\!\left[ \int_t^T \ell(s, X^\alpha_s, \alpha_s)\,ds + g(X^\alpha_T) \,\bigg|\, X^\alpha_t = x \right].
\]
By the dynamic programming principle, $V$ formally satisfies the Hamilton--Jacobi--Bellman equation
\[
  \partial_t V(t, x) + \inf_{a \in A}\Big\{ b(t, x, a) \cdot \nabla V + \tfrac12\, \mathrm{tr}(\sigma \sigma^{\top}\, \nabla^2 V) + \ell(t, x, a) \Big\} = 0, \quad V(T, x) = g(x).
\]
In the case where $\sigma$ does not depend on $a$ (so-called \emph{drift control}), the HJB equation is of the semilinear form \eqref{eq:semilinpde} with
\[
  f(t, x, y, z) = \inf_{a \in A}\!\Big\{ b(t, x, a) \cdot \sigma^{-\top}(t, x)\, z + \ell(t, x, a) \Big\}
\]
(assuming $\sigma$ is invertible; if not, work with $\sigma^{\top} \nabla V = Z$ as the primary object and a generalized Hamiltonian). The BSDE associated with this driver therefore gives a probabilistic representation of $V$ that bypasses the regularity demands of viscosity-solution theory, and produces an optimal feedback control $\alpha^*(s, X_s) \in \arg\min_a \{\cdots\}$ as a measurable selector.

When $\sigma$ \emph{does} depend on $a$, the BSDE picture must be replaced by a \emph{second-order} BSDE (2BSDE, Soner--Touzi--Zhang); this is one of the threads of \S7.

\subsection{Deep BSDE for high-dimensional PDEs}

The PDE \eqref{eq:semilinpde} in high spatial dimension $n$ is out of reach for grid-based solvers: storage is $\Theta(N^n)$ for $N$ points per axis. The Deep BSDE method of Han, Jentzen, and E (\cite{han-jentzen-e}) recasts solving the PDE as a single global stochastic optimization. Discretize time as $0 = t_0 < t_1 < \cdots < t_M = T$ with $\Delta t = T/M$. Parameterize:
\begin{itemize}
  \item the initial value by a scalar $\theta_0 \in \R$ (the candidate for $u(0, x_0)$);
  \item the spatial gradient $z(t, x) := \sigma^{\top}(t, x)\,\nabla u(t, x)$ by a neural network $z^\theta(t, x)$ with parameters $\theta$.
\end{itemize}
Simulate the forward Euler--Maruyama discretization of $X$ starting from $x_0$, and propagate
\[
  Y^\theta_{t_{k+1}} = Y^\theta_{t_k} - f\!\big(t_k, X_{t_k}, Y^\theta_{t_k}, z^\theta(t_k, X_{t_k})\big)\, \Delta t + z^\theta(t_k, X_{t_k}) \cdot \Delta B_k,
\]
with $Y^\theta_0 = \theta_0$. Train $\theta = (\theta_0, \theta_{\rm net})$ by stochastic gradient descent on the terminal-condition mismatch loss
\[
  \mathcal{J}(\theta) := \E\!\left[ \big| g(X_T) - Y^\theta_T \big|^2 \right].
\]
At a minimizer, $\theta_0 \approx u(0, x_0)$ and $z^\theta(\cdot, \cdot) \approx \sigma^{\top} \nabla u$.

The mathematical content is the \emph{variational characterization}: the function $u(0, x_0)$ is the unique value such that, when $Y_0 = u(0, x_0)$ and $Z$ is the BSDE solution, the residual $g(X_T) - Y_T$ vanishes a.s. Replacing ``a.s.'' by ``in $L^2$'' gives the loss, and replacing ``BSDE solution $Z$'' by ``best-NN approximation $z^\theta$'' gives the algorithm. The connection between $L^2$-loss and a.s.\ equality is precisely the BSDE uniqueness lemma in disguise.

In practice, Deep BSDE handles $n = 100$ on a laptop in minutes, well beyond the reach of finite differences. Reported accuracy is $10^{-3}$ to $10^{-4}$ relative error on benchmark problems (HJB, Black--Scholes with default risk, Allen--Cahn).

\begin{figure}[h]
\centering
\begin{tikzpicture}[
  >={Stealth},
  yk/.style={draw, fill=red!12, circle, minimum size=8mm, inner sep=0pt, font=\small},
  nn/.style={draw, fill=blue!12, rounded corners, minimum height=7mm, minimum width=9mm, font=\small},
  xk/.style={draw, fill=gray!12, rounded corners, minimum height=7mm, minimum width=9mm, font=\small},
]
% Y row
\node[yk] (Y0) at (0,0) {$Y_0$};
\node[nn] (NN0) at (1.6, 0) {$z^\theta_0$};
\node[yk] (Y1) at (3.2, 0) {$Y_1$};
\node[nn] (NN1) at (4.8, 0) {$z^\theta_1$};
\node[yk] (Y2) at (6.4, 0) {$Y_2$};
\node      (dots) at (7.7, 0) {$\cdots$};
\node[nn] (NNM)  at (9.0, 0) {$z^\theta_{M-1}$};
\node[yk] (YM)  at (10.6, 0) {$Y_M$};

\node[below=1pt of Y0, font=\scriptsize] {$=\theta_0$};

\draw[->] (Y0) -- (NN0);
\draw[->] (NN0) -- (Y1);
\draw[->] (Y1) -- (NN1);
\draw[->] (NN1) -- (Y2);
\draw[->] (Y2) -- (dots);
\draw[->] (dots) -- (NNM);
\draw[->] (NNM) -- (YM);

% X row
\node[xk] (X0) at (0,-1.8) {$X_0=x_0$};
\node[xk] (X1) at (3.2,-1.8) {$X_1$};
\node[xk] (X2) at (6.4,-1.8) {$X_2$};
\node     (Xd) at (7.7,-1.8) {$\cdots$};
\node[xk] (XM) at (10.6,-1.8) {$X_M$};

\draw[->, gray!70!black] (X0) -- node[font=\scriptsize, above]{$\Delta B_0$} (X1);
\draw[->, gray!70!black] (X1) -- node[font=\scriptsize, above]{$\Delta B_1$} (X2);
\draw[->, gray!70!black] (X2) -- (Xd);
\draw[->, gray!70!black] (Xd) -- (XM);

% Feed-in dashed arrows from X to NN
\draw[->, dashed, gray!70] (X0) -- (NN0);
\draw[->, dashed, gray!70] (X1) -- (NN1);
\draw[->, dashed, gray!70] (X2) -- (NNM);

% Loss
\node[draw, dashed, rounded corners, minimum width=4.0cm, minimum height=8mm, font=\small] (loss) at (8.0,-3.4) {$\mathcal{J}(\theta)=\E\big|g(X_M)-Y_M\big|^2$};
\draw[->] (XM) -- (loss.east);
\draw[->] (YM.south) to[bend left=18] (loss.north east);

% Legend
\node[font=\scriptsize, anchor=west] at (-0.7, 1.5) {trainable: $\theta_0\in\R$ (initial value), $\theta_{\rm net}$ (per-step nets)};
\node[font=\scriptsize, anchor=west] at (-0.7, -3.4) {input only: $\Delta B_k\sim\mathcal{N}(0,\Delta t)$};
\end{tikzpicture}
\caption{Architecture of the Deep BSDE method. The forward chain (gray) is the Euler--Maruyama discretization of $X$ driven by sampled Brownian increments $\Delta B_k$. The backward chain (red/blue) propagates $Y^\theta$ \emph{forward in time} from a trainable scalar $\theta_0$ using $Y^\theta_{k+1} = Y^\theta_k - f(t_k, X_k, Y^\theta_k, z^\theta_k)\Delta t + z^\theta_k \cdot \Delta B_k$, where each $z^\theta_k$ is a small neural network mapping $X_k \mapsto \sigma^\top(t_k, X_k)\nabla u(t_k, X_k)$. Training by SGD on the terminal mismatch loss $\mathcal{J}(\theta)$ drives $\theta_0 \to u(0, x_0)$. The variational characterization that makes this work is, in essence, the uniqueness half of Theorem~\ref{thm:pp}.}
\label{fig:deepbsde}
\end{figure}

\section{Outlook}

We close with four threads that the next article in this series could follow.

\paragraph{Forward--backward coupling (FBSDEs).} If the forward coefficients depend on $Y$ or $Z$, we obtain a \emph{coupled} forward--backward SDE:
\begin{align*}
  dX_s &= b(s, X_s, Y_s, Z_s)\,ds + \sigma(s, X_s, Y_s, Z_s)\, dB_s, \\
  -dY_s &= f(s, X_s, Y_s, Z_s)\,ds - Z_s\, dB_s, \qquad Y_T = g(X_T).
\end{align*}
This is the natural language of fully coupled stochastic optimal control via Pontryagin's maximum principle. Existence requires either a monotonicity condition (Ma--Yong), a small-time hypothesis (Antonelli), or the four-step scheme (Ma--Protter--Yong) when the PDE side has a classical solution.

\paragraph{Mean-field BSDEs.} Letting the driver depend on the marginal law $\mathcal{L}(Y_s)$ leads to McKean--Vlasov BSDEs, the natural language of mean-field games. The corresponding PDE lives on $[0, T] \times \R^n \times \mathcal{P}_2(\R^n)$ --- the so-called master equation --- and links back to the Wasserstein calculus of Lions.

\paragraph{Reflected BSDEs and American options.} A reflected BSDE (El~Karoui--Kapoudjian--Pardoux--Peng--Quenez 1997) adds an obstacle process $L = (L_t)$ to the equation: a solution is a triple $(Y, Z, K)$ where $K$ is a continuous nondecreasing process with $K_0 = 0$, $Y_t \ge L_t$ a.s., and
\[
  Y_t = \xi + \int_t^T f(s, Y_s, Z_s)\,ds + K_T - K_t - \int_t^T Z_s\, dB_s,
\]
together with the Skorokhod condition $\int_0^T (Y_t - L_t)\,dK_t = 0$. The process $K$ pushes $Y$ upward minimally, only when $Y$ touches the obstacle. The associated PDE is the obstacle problem
\[
  \min\!\big\{ -\partial_t u - \Lc u - f(t, x, u, \sigma^{\top}\nabla u),\ u - L \big\} = 0, \qquad u(T, x) = g(x),
\]
and the BSDE representation gives $Y_t = \mathrm{esssup}_{\tau \in [t, T]} \E[\int_t^\tau f\, ds + L_\tau \mathbf{1}_{\tau < T} + \xi \mathbf{1}_{\tau = T} \mid \F_t]$, i.e.\ the value of an optimal stopping problem with running cost $f$. The textbook application is American-option pricing: $L_t = (K - X_t)^+$ for a put with strike $K$, $f \equiv -r y$ for discounting; the early-exercise boundary is the contact set $\{Y = L\}$. Double-obstacle (Dynkin game) and second-order variants follow the same template.

\paragraph{Rough paths and rough BSDEs.} If the driving noise is rougher than Brownian motion (e.g.\ fractional Brownian motion with Hurst parameter $H < 1/2$), the It\^o calculus underlying BSDEs breaks down. Lyons' rough path theory and Hairer's regularity structures supply the replacement; rough BSDEs are an active area, with the Diehl--Friz--Gassiat formulation as the standard reference.

\paragraph{Acknowledgments.} We thank our reader for following the series. As always, comments, corrections, and counter-examples are welcome.

\begin{thebibliography}{9}
\bibitem{pp1990} Pardoux, \'E., Peng, S. \emph{Adapted solution of a backward stochastic differential equation.} Systems Control Lett.\ 14 (1990), 55--61.
\bibitem{el-karoui-quenez} El Karoui, N., Peng, S., Quenez, M.~C. \emph{Backward stochastic differential equations in finance.} Math.\ Finance 7 (1997), 1--71.
\bibitem{han-jentzen-e} Han, J., Jentzen, A., E, W. \emph{Solving high-dimensional partial differential equations using deep learning.} PNAS 115 (2018), 8505--8510.
\bibitem{pardoux-rascanu} Pardoux, \'E., R\u{a}\c{s}canu, A. \emph{Stochastic Differential Equations, Backward SDEs, Partial Differential Equations.} Springer, 2014.
\bibitem{revuz-yor} Revuz, D., Yor, M. \emph{Continuous Martingales and Brownian Motion}, 3rd ed., Springer, 1999.
\end{thebibliography}

\end{document}
